% GENERATED FILE. Edit canonical lessons, then run npm run generate:books.
% canonical-source-hash: fnv1a64:0d8ab9604fc333e0
% canonical-lessons: JA-C413-heiya, JA-C413-bonchi, JA-C413-hantou, JA-C413-tairiku, JA-C413-sabaku, JA-R413-review-408-413

\chapter{Plains, Peninsulas and Deserts}
\label{ch:ja-plains-peninsulas-and-deserts}
\hlchaptermodality{413}

\begin{chapteropening}
\textbf{By the end of this chapter:} \emph{I can describe the lie of the land.}

Complete the last lesson of chapter 413: I can describe the lie of the land.
\end{chapteropening}

\section[\texorpdfstring{\ja{へいや} (heiya)}{へいや (heiya)}]{\ja{へいや} (heiya) — a plain (flat land)}

\label{lesson:JA-C413-heiya}



\begin{quote}
Before the new one: say the Japanese for a plateau, then the Japanese for mountain climbing.
\end{quote}

\subsection*{You'll want to know: \ja{へいや}}
\textbf{\ja{へいや}} — \emph{heiya} — \textquotedblleft{}a plain (flat land)\textquotedblright{}.

\begin{quote}
\textbf{\ja{ひろい} \ja{へいやに} \ja{たんぼが} \ja{あります。}} — \emph{hiroi heiya ni tanbo ga arimasu} — \textquotedblleft{}There are rice fields on the wide plain.\textquotedblright{}
\end{quote}

\begin{cousinweb}
\textbf{\ja{へいや}} (\emph{heiya}) is Sino-Japanese (kango): two Chinese characters read with their on'yomi, readings borrowed from Chinese. They are \emph{hei}, \textquotedblleft{}flat\textquotedblright{} (Middle Chinese \emph{bjaeng}) and \emph{ya}, \textquotedblleft{}plains\textquotedblright{} (\emph{yaeX}) (KANJIDIC2; Middle Chinese in Baxter's transcription).
\end{cousinweb}

\begin{grammarlens}[title={What you've built}]
One more word for getting around.
\end{grammarlens}

\subsection*{Your turn}
\begin{itemize}
\raggedright
  \item \emph{Say it:} \emph{heiya}
  \item \emph{Say it:} \emph{heiya}, once more
  \item \emph{Say it:} \emph{hiroi heiya ni tanbo ga arimasu} — There are rice fields on the wide plain.
  \item \emph{Recall:} say the Japanese for a plateau, then the Japanese for mountain climbing, then say \emph{heiya} again
\end{itemize}

\begin{tcolorbox}[breakable,colback=teal!4,colframe=teal!35!black,title={Before you move on}]
What does \ja{へいや} mean? (\textquotedblleft{}A plain (flat land)\textquotedblright{}.) Say it once more.
\end{tcolorbox}
\section[\texorpdfstring{\ja{ぼんち} (bonchi)}{ぼんち (bonchi)}]{\ja{ぼんち} (bonchi) — a basin among mountains}

\label{lesson:JA-C413-bonchi}



\begin{quote}
Before the new one: say the Japanese for mountain climbing, then the Japanese for a plain.
\end{quote}

\subsection*{You'll want to know: \ja{ぼんち}}
\textbf{\ja{ぼんち}} — \emph{bonchi} — \textquotedblleft{}a basin among mountains\textquotedblright{}.

\begin{quote}
\textbf{\ja{ぼんちの} \ja{なつは} \ja{むしあついです。}} — \emph{bonchi no natsu wa mushiatsui desu} — \textquotedblleft{}Summers in the basin are hot and humid.\textquotedblright{}
\end{quote}

\begin{cousinweb}
\textbf{\ja{ぼんち}} (\emph{bonchi}) is Sino-Japanese (kango): two Chinese characters read with their on'yomi, readings borrowed from Chinese. They are \emph{bon}, \textquotedblleft{}basin\textquotedblright{} (Middle Chinese \emph{bwon}) and \emph{chi}, \textquotedblleft{}ground\textquotedblright{} (\emph{dijH}) (KANJIDIC2; Middle Chinese in Baxter's transcription).
\end{cousinweb}

\begin{grammarlens}[title={What you've built}]
One more word for getting around.
\end{grammarlens}

\subsection*{Your turn}
\begin{itemize}
\raggedright
  \item \emph{Say it:} \emph{bonchi}
  \item \emph{Say it:} \emph{bonchi}, once more
  \item \emph{Say it:} \emph{bonchi no natsu wa mushiatsui desu} — Summers in the basin are hot and humid.
  \item \emph{Recall:} say the Japanese for mountain climbing, then the Japanese for a plain, then say \emph{bonchi} again
\end{itemize}

\begin{tcolorbox}[breakable,colback=teal!4,colframe=teal!35!black,title={Before you move on}]
What does \ja{ぼんち} mean? (\textquotedblleft{}A basin among mountains\textquotedblright{}.) Say it once more.
\end{tcolorbox}
\section[\texorpdfstring{\ja{はんとう} (hantō)}{はんとう (hantō)}]{\ja{はんとう} (hantō) — a peninsula}

\label{lesson:JA-C413-hantou}



\begin{quote}
Before the new one: say the Japanese for a plain, then the Japanese for a basin among mountains.
\end{quote}

\subsection*{You'll want to know: \ja{はんとう}}
\textbf{\ja{はんとう}} — \emph{hantō} — \textquotedblleft{}a peninsula\textquotedblright{}.

\begin{quote}
\textbf{\ja{この} \ja{はんとうには} \ja{みさきが} \ja{みっつ} \ja{あります。}} — \emph{kono hantō ni wa misaki ga mittsu arimasu} — \textquotedblleft{}This peninsula has three capes.\textquotedblright{}
\end{quote}

\begin{cousinweb}
\textbf{\ja{はんとう}} (\emph{hantō}) is Sino-Japanese (kango): two Chinese characters read with their on'yomi, readings borrowed from Chinese. They are \emph{han}, \textquotedblleft{}half\textquotedblright{} (Middle Chinese \emph{panH}) and \emph{tō}, \textquotedblleft{}island\textquotedblright{} (\emph{tawX}) (KANJIDIC2; Middle Chinese in Baxter's transcription).
\end{cousinweb}

\begin{grammarlens}[title={What you've built}]
One more word for getting around.
\end{grammarlens}

\subsection*{Your turn}
\begin{itemize}
\raggedright
  \item \emph{Say it:} \emph{hantō}
  \item \emph{Say it:} \emph{hantō}, once more
  \item \emph{Say it:} \emph{kono hantō ni wa misaki ga mittsu arimasu} — This peninsula has three capes.
  \item \emph{Recall:} say the Japanese for a plain, then the Japanese for a basin among mountains, then say \emph{hantō} again
\end{itemize}

\begin{tcolorbox}[breakable,colback=teal!4,colframe=teal!35!black,title={Before you move on}]
What does \ja{はんとう} mean? (\textquotedblleft{}A peninsula\textquotedblright{}.) Say it once more.
\end{tcolorbox}
\section[\texorpdfstring{\ja{たいりく} (tairiku)}{たいりく (tairiku)}]{\ja{たいりく} (tairiku) — a continent}

\label{lesson:JA-C413-tairiku}



\begin{quote}
Before the new one: say the Japanese for a basin among mountains, then the Japanese for a peninsula.
\end{quote}

\subsection*{You'll want to know: \ja{たいりく}}
\textbf{\ja{たいりく}} — \emph{tairiku} — \textquotedblleft{}a continent\textquotedblright{}.

\begin{quote}
\textbf{\ja{ふねで} \ja{たいりくに} \ja{わたりました。}} — \emph{fune de tairiku ni watarimashita} — \textquotedblleft{}We crossed over to the continent by ship.\textquotedblright{}
\end{quote}

\begin{cousinweb}
\textbf{\ja{たいりく}} (\emph{tairiku}) is Sino-Japanese (kango): two Chinese characters read with their on'yomi, readings borrowed from Chinese. They are \emph{tai}, \textquotedblleft{}large\textquotedblright{} (Middle Chinese \emph{dajH}) and \emph{riku}, \textquotedblleft{}land\textquotedblright{} (\emph{ljuwk}) (KANJIDIC2; Middle Chinese in Baxter's transcription).
\end{cousinweb}

\begin{grammarlens}[title={What you've built}]
One more word for getting around.
\end{grammarlens}

\subsection*{Your turn}
\begin{itemize}
\raggedright
  \item \emph{Say it:} \emph{tairiku}
  \item \emph{Say it:} \emph{tairiku}, once more
  \item \emph{Say it:} \emph{fune de tairiku ni watarimashita} — We crossed over to the continent by ship.
  \item \emph{Recall:} say the Japanese for a basin among mountains, then the Japanese for a peninsula, then say \emph{tairiku} again
\end{itemize}

\begin{tcolorbox}[breakable,colback=teal!4,colframe=teal!35!black,title={Before you move on}]
What does \ja{たいりく} mean? (\textquotedblleft{}A continent\textquotedblright{}.) Say it once more.
\end{tcolorbox}
\section[\texorpdfstring{\ja{さばく} (sabaku)}{さばく (sabaku)}]{\ja{さばく} (sabaku) — a desert}

\label{lesson:JA-C413-sabaku}



\begin{quote}
Before the new one: say the Japanese for a peninsula, then the Japanese for a continent.
\end{quote}

\subsection*{You'll want to know: \ja{さばく}}
\textbf{\ja{さばく}} — \emph{sabaku} — \textquotedblleft{}a desert\textquotedblright{}.

\begin{quote}
\textbf{\ja{さばくは} \ja{ひるは} \ja{あついです。}} — \emph{sabaku wa hiru wa atsui desu} — \textquotedblleft{}The desert is hot during the day.\textquotedblright{}
\end{quote}

\begin{cousinweb}
\textbf{\ja{さばく}} (\emph{sabaku}) is Sino-Japanese (kango): two Chinese characters read with their on'yomi, readings borrowed from Chinese. They are \emph{sa}, \textquotedblleft{}sand\textquotedblright{} (Middle Chinese \emph{srae}) and \emph{baku}, \textquotedblleft{}desert\textquotedblright{} (\emph{mak}) (KANJIDIC2; Middle Chinese in Baxter's transcription).
\end{cousinweb}

\begin{grammarlens}[title={What you've built}]
One more word for getting around.
\end{grammarlens}

\subsection*{Your turn}
\begin{itemize}
\raggedright
  \item \emph{Say it:} \emph{sabaku}
  \item \emph{Say it:} \emph{sabaku}, once more
  \item \emph{Say it:} \emph{sabaku wa hiru wa atsui desu} — The desert is hot during the day.
  \item \emph{Recall:} say the Japanese for a peninsula, then the Japanese for a continent, then say \emph{sabaku} again
\end{itemize}

\begin{tcolorbox}[breakable,colback=teal!4,colframe=teal!35!black,title={Before you move on}]
What does \ja{さばく} mean? (\textquotedblleft{}A desert\textquotedblright{}.) Say it once more.
\end{tcolorbox}
\section[(dialogue)]{Looking back: from Incense and Candles to Plains, Peninsulas and Deserts}

\label{lesson:JA-R413-review-408-413}



\begin{quote}
Say the Japanese for an incense stick and a candle.
\end{quote}

\subsection*{Your turn}
Say these aloud:

\begin{itemize}
\raggedright
  \item an incense stick, a candle, a household altar, religion, to enshrine
  \item an area, a capital city, a national border, a boundary, a dialect
  \item one's home area, one's place of origin, a field (for crops), a meadow, to grow thick
  \item a stream, a spring (of water), to become muddy, to surround, to scoop up
  \item a mountain pass, a summit, a mountain range, a plateau, mountain climbing
  \item a plain, a basin among mountains, a peninsula, a continent, a desert
\end{itemize}

\begin{tcolorbox}[breakable,colback=teal!4,colframe=teal!35!black,title={Before you move on}]
An incense stick? (\textbf{\ja{せんこう}}, \emph{senkō}.) A candle? (\textbf{\ja{ろうそく}}, \emph{rōsoku}.) A household altar? (\textbf{\ja{ぶつだん}}, \emph{butsudan}.) Religion? (\textbf{\ja{しゅうきょう}}, \emph{shūkyō}.) To enshrine? (\textbf{\ja{まつる}}, \emph{matsuru}.) An area? (\textbf{\ja{ちいき}}, \emph{chiiki}.) A capital city? (\textbf{\ja{しゅと}}, \emph{shuto}.) A national border? (\textbf{\ja{こっきょう}}, \emph{kokkyō}.) A boundary? (\textbf{\ja{さかい}}, \emph{sakai}.) A dialect? (\textbf{\ja{ほうげん}}, \emph{hōgen}.) One's home area? (\textbf{\ja{じもと}}, \emph{jimoto}.) One's place of origin? (\textbf{\ja{しゅっしん}}, \emph{shusshin}.) A field (for crops)? (\textbf{\ja{はたけ}}, \emph{hatake}.) A meadow? (\textbf{\ja{のはら}}, \emph{nohara}.) To grow thick? (\textbf{\ja{しげる}}, \emph{shigeru}.) A stream? (\textbf{\ja{おがわ}}, \emph{ogawa}.) A spring (of water)? (\textbf{\ja{いずみ}}, \emph{izumi}.) To become muddy? (\textbf{\ja{にごる}}, \emph{nigoru}.) To surround? (\textbf{\ja{かこむ}}, \emph{kakomu}.) To scoop up? (\textbf{\ja{すくう}}, \emph{sukuu}.) A mountain pass? (\textbf{\ja{とうげ}}, \emph{tōge}.) A summit? (\textbf{\ja{ちょうじょう}}, \emph{chōjō}.) A mountain range? (\textbf{\ja{さんみゃく}}, \emph{sanmyaku}.) A plateau? (\textbf{\ja{こうげん}}, \emph{kōgen}.) Mountain climbing? (\textbf{\ja{とざん}}, \emph{tozan}.) A plain? (\textbf{\ja{へいや}}, \emph{heiya}.) A basin among mountains? (\textbf{\ja{ぼんち}}, \emph{bonchi}.) A peninsula? (\textbf{\ja{はんとう}}, \emph{hantō}.) A continent? (\textbf{\ja{たいりく}}, \emph{tairiku}.) A desert? (\textbf{\ja{さばく}}, \emph{sabaku}.)
\end{tcolorbox}
